% Options for packages loaded elsewhere
\PassOptionsToPackage{unicode}{hyperref}
\PassOptionsToPackage{hyphens}{url}
\PassOptionsToPackage{dvipsnames,svgnames,x11names}{xcolor}
%
\documentclass[
]{article}
\usepackage{amsmath,amssymb}
\usepackage{iftex}
\ifPDFTeX
  \usepackage[T1]{fontenc}
  \usepackage[utf8]{inputenc}
  \usepackage{textcomp} % provide euro and other symbols
\else % if luatex or xetex
  \usepackage{unicode-math} % this also loads fontspec
  \defaultfontfeatures{Scale=MatchLowercase}
  \defaultfontfeatures[\rmfamily]{Ligatures=TeX,Scale=1}
\fi
\usepackage{lmodern}
\ifPDFTeX\else
  % xetex/luatex font selection
  \setmainfont[]{DejaVu Serif}
  \setsansfont[]{DejaVu Sans}
  \setmonofont[]{DejaVu Sans Mono}
\fi
% Use upquote if available, for straight quotes in verbatim environments
\IfFileExists{upquote.sty}{\usepackage{upquote}}{}
\IfFileExists{microtype.sty}{% use microtype if available
  \usepackage[]{microtype}
  \UseMicrotypeSet[protrusion]{basicmath} % disable protrusion for tt fonts
}{}
\makeatletter
\@ifundefined{KOMAClassName}{% if non-KOMA class
  \IfFileExists{parskip.sty}{%
    \usepackage{parskip}
  }{% else
    \setlength{\parindent}{0pt}
    \setlength{\parskip}{6pt plus 2pt minus 1pt}}
}{% if KOMA class
  \KOMAoptions{parskip=half}}
\makeatother
\usepackage{xcolor}
\usepackage[margin=0.7in]{geometry}
\usepackage{color}
\usepackage{fancyvrb}
\newcommand{\VerbBar}{|}
\newcommand{\VERB}{\Verb[commandchars=\\\{\}]}
\DefineVerbatimEnvironment{Highlighting}{Verbatim}{commandchars=\\\{\}}
% Add ',fontsize=\small' for more characters per line
\usepackage{framed}
\definecolor{shadecolor}{RGB}{248,248,248}
\newenvironment{Shaded}{\begin{snugshade}}{\end{snugshade}}
\newcommand{\AlertTok}[1]{\textcolor[rgb]{0.94,0.16,0.16}{#1}}
\newcommand{\AnnotationTok}[1]{\textcolor[rgb]{0.56,0.35,0.01}{\textbf{\textit{#1}}}}
\newcommand{\AttributeTok}[1]{\textcolor[rgb]{0.13,0.29,0.53}{#1}}
\newcommand{\BaseNTok}[1]{\textcolor[rgb]{0.00,0.00,0.81}{#1}}
\newcommand{\BuiltInTok}[1]{#1}
\newcommand{\CharTok}[1]{\textcolor[rgb]{0.31,0.60,0.02}{#1}}
\newcommand{\CommentTok}[1]{\textcolor[rgb]{0.56,0.35,0.01}{\textit{#1}}}
\newcommand{\CommentVarTok}[1]{\textcolor[rgb]{0.56,0.35,0.01}{\textbf{\textit{#1}}}}
\newcommand{\ConstantTok}[1]{\textcolor[rgb]{0.56,0.35,0.01}{#1}}
\newcommand{\ControlFlowTok}[1]{\textcolor[rgb]{0.13,0.29,0.53}{\textbf{#1}}}
\newcommand{\DataTypeTok}[1]{\textcolor[rgb]{0.13,0.29,0.53}{#1}}
\newcommand{\DecValTok}[1]{\textcolor[rgb]{0.00,0.00,0.81}{#1}}
\newcommand{\DocumentationTok}[1]{\textcolor[rgb]{0.56,0.35,0.01}{\textbf{\textit{#1}}}}
\newcommand{\ErrorTok}[1]{\textcolor[rgb]{0.64,0.00,0.00}{\textbf{#1}}}
\newcommand{\ExtensionTok}[1]{#1}
\newcommand{\FloatTok}[1]{\textcolor[rgb]{0.00,0.00,0.81}{#1}}
\newcommand{\FunctionTok}[1]{\textcolor[rgb]{0.13,0.29,0.53}{\textbf{#1}}}
\newcommand{\ImportTok}[1]{#1}
\newcommand{\InformationTok}[1]{\textcolor[rgb]{0.56,0.35,0.01}{\textbf{\textit{#1}}}}
\newcommand{\KeywordTok}[1]{\textcolor[rgb]{0.13,0.29,0.53}{\textbf{#1}}}
\newcommand{\NormalTok}[1]{#1}
\newcommand{\OperatorTok}[1]{\textcolor[rgb]{0.81,0.36,0.00}{\textbf{#1}}}
\newcommand{\OtherTok}[1]{\textcolor[rgb]{0.56,0.35,0.01}{#1}}
\newcommand{\PreprocessorTok}[1]{\textcolor[rgb]{0.56,0.35,0.01}{\textit{#1}}}
\newcommand{\RegionMarkerTok}[1]{#1}
\newcommand{\SpecialCharTok}[1]{\textcolor[rgb]{0.81,0.36,0.00}{\textbf{#1}}}
\newcommand{\SpecialStringTok}[1]{\textcolor[rgb]{0.31,0.60,0.02}{#1}}
\newcommand{\StringTok}[1]{\textcolor[rgb]{0.31,0.60,0.02}{#1}}
\newcommand{\VariableTok}[1]{\textcolor[rgb]{0.00,0.00,0.00}{#1}}
\newcommand{\VerbatimStringTok}[1]{\textcolor[rgb]{0.31,0.60,0.02}{#1}}
\newcommand{\WarningTok}[1]{\textcolor[rgb]{0.56,0.35,0.01}{\textbf{\textit{#1}}}}
\usepackage{longtable,booktabs,array}
\usepackage{calc} % for calculating minipage widths
% Correct order of tables after \paragraph or \subparagraph
\usepackage{etoolbox}
\makeatletter
\patchcmd\longtable{\par}{\if@noskipsec\mbox{}\fi\par}{}{}
\makeatother
% Allow footnotes in longtable head/foot
\IfFileExists{footnotehyper.sty}{\usepackage{footnotehyper}}{\usepackage{footnote}}
\makesavenoteenv{longtable}
\setlength{\emergencystretch}{3em} % prevent overfull lines
\providecommand{\tightlist}{%
  \setlength{\itemsep}{0pt}\setlength{\parskip}{0pt}}
\setcounter{secnumdepth}{-\maxdimen} % remove section numbering
\ifLuaTeX
\usepackage[bidi=basic]{babel}
\else
\usepackage[bidi=default]{babel}
\fi
\babelprovide[main,import]{russian}
\ifPDFTeX
\else
\babelfont{rm}[]{DejaVu Serif}
\fi
% get rid of language-specific shorthands (see #6817):
\let\LanguageShortHands\languageshorthands
\def\languageshorthands#1{}
\ifLuaTeX
  \usepackage{selnolig}  % disable illegal ligatures
\fi
\usepackage{bookmark}
\IfFileExists{xurl.sty}{\usepackage{xurl}}{} % add URL line breaks if available
\urlstyle{same}
\hypersetup{
  pdflang={ru-RU},
  colorlinks=true,
  linkcolor={blue},
  filecolor={Maroon},
  citecolor={Blue},
  urlcolor={blue},
  pdfcreator={LaTeX via pandoc}}

\author{}
\date{}

\begin{document}

\section{1. Общая
идея}\label{ux43eux431ux449ux430ux44f-ux438ux434ux435ux44f}

Программа делится на независимые модули, каждый отвечает за одну задачу.
Задание требует:

\begin{itemize}
\tightlist
\item
  модуль управления ходом моделирования;
\item
  модуль контроля соблюдения правил;
\item
  модуль риэлторского агентства;
\item
  несколько модулей-игроков с разными стратегиями.
\end{itemize}

Кроме них есть настройки игры, состояние игры и интерфейс.

\section{2. Что требует задание и где это
реализовано}\label{ux447ux442ux43e-ux442ux440ux435ux431ux443ux435ux442-ux437ux430ux434ux430ux43dux438ux435-ux438-ux433ux434ux435-ux44dux442ux43e-ux440ux435ux430ux43bux438ux437ux43eux432ux430ux43dux43e}

\begin{longtable}[]{@{}
  >{\raggedright\arraybackslash}p{(\columnwidth - 2\tabcolsep) * \real{0.5000}}
  >{\raggedright\arraybackslash}p{(\columnwidth - 2\tabcolsep) * \real{0.5000}}@{}}
\toprule\noalign{}
\begin{minipage}[b]{\linewidth}\raggedright
Требование задания
\end{minipage} & \begin{minipage}[b]{\linewidth}\raggedright
Где реализовано
\end{minipage} \\
\midrule\noalign{}
\endhead
\bottomrule\noalign{}
\endlastfoot
Игра на 2--5 игроков, 6--24 месяца & \texttt{Config}, проверка при
создании игры в \texttt{GameEngine} \\
В начале игры игроки решают, куда вложить деньги & начальная фаза
(раздел 5.1) \\
Игроки-люди и игроки-программы & \texttt{Player}: \texttt{HumanPlayer} и
программные стратегии \\
Управление ходом моделирования & \texttt{GameEngine} \\
Контроль соблюдения правил & \texttt{RulesValidator} \\
Риэлторское агентство & \texttt{RealtorAgency} \\
Несколько стратегий игры & \texttt{BuilderPlayer},
\texttt{MarketPlayer}, \texttt{BalancedPlayer} \\
Разные типы домов и супермаркеты, изменяемые стоимость и сроки &
\texttt{Building}, \texttt{House}, \texttt{Supermarket},
\texttt{Config} \\
Равномерные затраты на стройку & \texttt{GameEngine}, шаг 1 месяца \\
Сезонный спрос и продажи, влияние супермаркетов и домов &
\texttt{District}, формулы \texttt{math-model.md} (разделы 3 и 9) \\
Реклама жилья и супермаркетов & \texttt{Decision}, \texttt{GameEngine},
\texttt{RealtorAgency} \\
Подсчёт капитала и победитель & \texttt{GameEngine}, раздел 6 \\
Информация о чужих объектах и продажах & \texttt{PublicView} \\
\end{longtable}

\section{3. Модули}\label{ux43cux43eux434ux443ux43bux438}

\subsection{3.1. GameEngine --- координатор
моделирования}\label{gameengine-ux43aux43eux43eux440ux434ux438ux43dux430ux442ux43eux440-ux43cux43eux434ux435ux43bux438ux440ux43eux432ux430ux43dux438ux44f}

\textbf{Задача:} координировать ход игры. Сам он не содержит всей
математики: он вызывает игроков, \texttt{RulesValidator} и
\texttt{RealtorAgency} в нужном порядке и применяет их результаты к
состоянию.

\begin{itemize}
\tightlist
\item
  создаёт игру по настройкам (\(N\), \(M\), стартовый капитал) и
  проверяет допустимость \(N\) и \(M\);
\item
  проводит начальную фазу и затем каждый месяц (раздел 5);
\item
  списывает ежемесячные платежи за стройку, продвигает стройки,
  начисляет выручку и прибыль;
\item
  \textbf{замораживает стройку на месяц}, если на её платёж не хватает
  денег (платёж не списывается, срок сдвигается на месяц). Это
  автоматическое следствие, а не проверка решения игрока;
\item
  считает капитал (раздел 6) и после месяца \(M\) определяет победителя.
\end{itemize}

\subsection{3.2. RulesValidator --- контроль
правил}\label{rulesvalidator-ux43aux43eux43dux442ux440ux43eux43bux44c-ux43fux440ux430ux432ux438ux43b}

\textbf{Задача:} проверять \textbf{решения игроков} до того, как они
попадут в игру.

\begin{longtable}[]{@{}
  >{\raggedright\arraybackslash}p{(\columnwidth - 2\tabcolsep) * \real{0.5000}}
  >{\raggedright\arraybackslash}p{(\columnwidth - 2\tabcolsep) * \real{0.5000}}@{}}
\toprule\noalign{}
\begin{minipage}[b]{\linewidth}\raggedright
Проверка
\end{minipage} & \begin{minipage}[b]{\linewidth}\raggedright
Что происходит при нарушении
\end{minipage} \\
\midrule\noalign{}
\endhead
\bottomrule\noalign{}
\endlastfoot
Денег на рекламу не хватает & реклама уменьшается до допустимой или
отклоняется \\
Денег на платежи за новые стройки в следующем месяце не хватает & новая
стройка не принимается \\
Выставлено на продажу больше, чем доступно & объём обрезается до
доступного \\
Цена или объём не положительны & предложение отклоняется \\
\textbf{Цена выше предельной} & предложение отклоняется сразу, игроку
сообщается причина \\
Реклама супермаркетов без построенного супермаркета & реклама
отклоняется \\
Решение подано не вовремя или повторно & отклоняется \\
\end{longtable}

Как проверяется стройка. Первый платёж за новую стройку списывается в
шаге 1 \textbf{следующего} месяца, до поступления выручки. Поэтому
валидатор проверяет, что денег, оставшихся после рекламы этого месяца,
хватит на все платежи за стройки в следующем месяце (старые и новые).
Недостаток денег на \textbf{уже идущую} стройку --- не нарушение правил
игроком, это случай заморозки (см. \texttt{GameEngine}).

\subsection{3.3. RealtorAgency --- риэлторское
агентство}\label{realtoragency-ux440ux438ux44dux43bux442ux43eux440ux441ux43aux43eux435-ux430ux433ux435ux43dux442ux441ux442ux432ux43e}

\textbf{Задача:} распределить спрос между предложениями игроков по
правилам \texttt{math-model.md} (раздел 6, рис. 4).

\textbf{Спрос по типам домов.} Для каждого типа дома (панельный,
монолитный, кирпичный) спрос считается \textbf{отдельно}
(\texttt{math-model.md}, раздел 3): общий спрос на все типы вместе в три
раза больше базового. Агентство обрабатывает каждый тип независимо.

Алгоритм для одного типа дома:

\begin{enumerate}
\def\labelenumi{\arabic{enumi}.}
\tightlist
\item
  взять спрос \(D\) на этот тип и список предложений (игрок, количество
  квартир, цена за кв. м, рекламный бонус);
\item
  отсортировать предложения по цене от низкой к высокой;
\item
  по очереди удовлетворять спрос: каждый игрок получает
  \(\min(\text{его предложение}, \text{остаток спроса})\); при равных
  ценах остаток делится пропорционально предложению;
\item
  \textbf{реклама.} После распределения проданное количество каждого
  игрока умножается на \((1 + b)\), где \(b\) --- рекламный бонус
  игрока, но не больше его предложения. Дополнительные продажи за счёт
  рекламы \textbf{не отнимаются у других игроков}: допущение модели ---
  реклама привлекает покупателей сверх базового спроса. Поэтому
  суммарные продажи могут немного превысить \(D\);
\item
  вернуть количество проданных квартир каждого игрока.
\end{enumerate}

Агентство ничего не знает ни о деньгах, ни о правилах: получает данные и
возвращает результат, поэтому его легко проверять отдельно. Предложения
с недопустимой ценой до него не доходят (их отклоняет
\texttt{RulesValidator}).

\subsection{3.4. Player ---
игроки}\label{player-ux438ux433ux440ux43eux43aux438}

\textbf{Задача:} в начале месяца принять решения. \texttt{Player} ---
базовый класс: он хранит данные игрока (имя, деньги, объекты, рекламный
бонус прошлого месяца) и объявляет метод принятия решения. Конкретные
стратегии наследуют его и реализуют метод по-своему. Все игроки
(программные и человек) получают \texttt{PublicView} и возвращают
\texttt{Decision}.

\texttt{Decision} содержит: - какие объекты начать строить (тип и
количество); - сколько квартир выставить на продажу и по какой цене за
кв. м (по типам домов); - сколько потратить на рекламу жилья; - сколько
потратить на рекламу товаров в супермаркетах (\textbf{общий бюджет
игрока} на все его супермаркеты: надбавка к прибыли применяется к
каждому его супермаркету, бюджет не делится).

\subsection{3.5. Config ---
настройки}\label{config-ux43dux430ux441ux442ux440ux43eux439ux43aux438}

Все числа хранятся в одном месте:

\begin{itemize}
\tightlist
\item
  \(N\), \(M\), стартовый капитал, месяц старта;
\item
  стоимость, срок, число квартир и площадь по типам домов; стоимость и
  срок супермаркета;
\item
  сезонные коэффициенты спроса и продаж, базовый спрос (на каждый тип
  дома) и базовая прибыль \textbf{одного} супермаркета (фиксированная
  сумма);
\item
  коэффициенты влияния (рост спроса от супермаркета, рост продаж от
  дома, вес строящегося объекта);
\item
  предельная цена и «справедливая» цена; правило ничьей;
\item
  расширение: включение задержек строительства и их вероятность.
\end{itemize}

\subsection{3.6. GameState и
PublicView}\label{gamestate-ux438-publicview}

\texttt{GameState} хранит всё состояние: игроков, объекты, историю
продаж и спроса.

\textbf{Одновременность.} Все игроки получают одинаковый
\texttt{PublicView}, составленный на начало месяца, и принимают решения
независимо друг от друга. Решения обрабатываются только после того, как
собраны от всех. Порядок игроков не влияет на результат.

\texttt{PublicView} содержит: - текущий месяц и календарный месяц; -
свои деньги, объекты и рекламный бонус; - по каждому другому игроку:
объекты (тип, готов или строится, сколько месяцев строится из \(T\)); -
историю продаж по всем игрокам за все прошедшие месяцы: игрок, тип дома,
количество проданных квартир, цена за кв. м; - историю спроса по типам
домов.

Деньги соперников не показываются.

\subsection{3.7.
Интерфейс}\label{ux438ux43dux442ux435ux440ux444ux435ux439ux441}

Показывает игроку-человеку \texttt{PublicView} и принимает его решение;
в конце показывает таблицу итогов (рис. 5). Для программных игроков
интерфейс не нужен: они работают напрямую, что удобно для быстрых
автоматических партий.

\textbf{Если решение человека нарушает правила.} Валидатор возвращает
причину, игроку даётся до 3 попыток исправить. После этого недопустимые
части решения отбрасываются (допустимые остаются, объёмы обрезаются до
допустимых).

\section{4. Данные}\label{ux434ux430ux43dux43dux44bux435}

\begin{longtable}[]{@{}
  >{\raggedright\arraybackslash}p{(\columnwidth - 2\tabcolsep) * \real{0.5000}}
  >{\raggedright\arraybackslash}p{(\columnwidth - 2\tabcolsep) * \real{0.5000}}@{}}
\toprule\noalign{}
\begin{minipage}[b]{\linewidth}\raggedright
Класс
\end{minipage} & \begin{minipage}[b]{\linewidth}\raggedright
Что хранит
\end{minipage} \\
\midrule\noalign{}
\endhead
\bottomrule\noalign{}
\endlastfoot
\texttt{Building} & тип, стоимость, срок, сколько месяцев строится,
владелец \\
\texttt{House} & как \texttt{Building}, плюс число квартир, площадь
квартиры, сколько продано \\
\texttt{Supermarket} & как \texttt{Building}, плюс признак «построен» \\
\texttt{Offer} & игрок, тип дома, количество квартир, цена за кв. м,
рекламный бонус \\
\texttt{Decision} & решения игрока за месяц \\
\texttt{District} & собирает данные по \textbf{всем} игрокам: число
готовых супермаркетов, строящихся супермаркетов, готовых домов,
строящихся домов. Нужен для формул спроса и прибыли супермаркетов (\(S\)
и \(H\) в \texttt{math-model.md}) \\
\texttt{Config} & настройки \\
\end{longtable}

Данные игрока (имя, деньги, объекты) хранятся в базовом классе
\texttt{Player}. Диаграмма классов --- на рис. 2.

\section{5. Ход игры}\label{ux445ux43eux434-ux438ux433ux440ux44b}

\subsection{5.1. Начальная фаза (до месяца
1)}\label{ux43dux430ux447ux430ux43bux44cux43dux430ux44f-ux444ux430ux437ux430-ux434ux43e-ux43cux435ux441ux44fux446ux430-1}

\begin{enumerate}
\def\labelenumi{\arabic{enumi}.}
\tightlist
\item
  каждый игрок получает стартовый капитал;
\item
  каждый игрок получает \texttt{PublicView} и возвращает
  \texttt{Decision}, в котором учитываются только \textbf{стартовые
  стройки} (продажи и реклама ещё невозможны);
\item
  \texttt{RulesValidator} проверяет, что денег хватит на платежи за эти
  стройки в месяце 1;
\item
  стройки создаются; их первые платежи списываются в шаге 1 месяца 1.
\end{enumerate}

\subsection{5.2. Один
месяц}\label{ux43eux434ux438ux43d-ux43cux435ux441ux44fux446}

\texttt{GameEngine} выполняет шаги в том же порядке, что и в
\texttt{math-model.md} (раздел 13, рис. 3):

\begin{enumerate}
\def\labelenumi{\arabic{enumi}.}
\tightlist
\item
  списываются затраты на строительство (по всем идущим стройкам;
  нехватка денег замораживает стройку);
\item
  зачисляется выручка от продаж жилья прошлого месяца;
\item
  зачисляется прибыль супермаркетов;
\item
  игроки получают \texttt{PublicView} и \textbf{одновременно} возвращают
  \texttt{Decision}: новые стройки, предложение жилья, цены, реклама;
\item
  \texttt{RulesValidator} проверяет решения и исправляет недопустимое;
\item
  после проверки снимаются деньги за рекламу жилья и супермаркетов
  (поэтому баланс не уходит в минус);
\item
  \texttt{RealtorAgency} распределяет спрос и возвращает продажи;
\item
  обновляются объекты (прогресс стройки, проданные квартиры) и капитал
  игроков.
\end{enumerate}

\textbf{Новые стройки.} Решение начать стройку принимается в шаге 4
месяца \(t\). Её первый платёж списывается в шаге 1 месяца \(t+1\) и с
этого месяца она продвигается.

\subsection{5.3.
Псевдокод}\label{ux43fux441ux435ux432ux434ux43eux43aux43eux434}

\begin{Shaded}
\begin{Highlighting}[]
\NormalTok{Setup:}
\NormalTok{    для каждого игрока: Decision = игрок.decide(PublicView)   \# только стройки}
\NormalTok{    проверить, создать стройки}

\NormalTok{для t = 1..M:}
\NormalTok{    списать платежи за стройки (нет денег {-}\textgreater{} заморозить стройку)}
\NormalTok{    зачислить выручку прошлого месяца}
\NormalTok{    зачислить прибыль супермаркетов}
\NormalTok{    view = PublicView(на начало месяца)}
\NormalTok{    для каждого игрока: Decision[i] = игрок[i].decide(view)     \# независимо}
\NormalTok{    для каждого игрока: Decision[i] = RulesValidator.check(Decision[i])}
\NormalTok{    списать рекламу}
\NormalTok{    для каждого типа дома j:}
\NormalTok{        продажи[j] = RealtorAgency.distribute(спрос[j], предложения[j])}
\NormalTok{    обновить объекты, продажи, капитал}

\NormalTok{зачислить выручку месяца M}
\NormalTok{победитель = игрок с максимальным капиталом (ничья возможна)}
\end{Highlighting}
\end{Shaded}

\section{6. Расчёт
капитала}\label{ux440ux430ux441ux447ux451ux442-ux43aux430ux43fux438ux442ux430ux43bux430}

Капитал считается по \texttt{math-model.md} (разделы 10--11):
\(K = C_{\text{construction}} + C_{\text{housing}} + C_{\text{markets}} + Cash\).

\begin{itemize}
\tightlist
\item
  \(C_{\text{construction}}\) --- только \textbf{незавершённые} объекты
  (и дома, и супермаркеты): потраченная на них сумма;
\item
  \(C_{\text{housing}}\) --- только \textbf{непроданное жильё в
  построенных домах} по себестоимости кв. м;
\item
  \(C_{\text{markets}}\) --- только \textbf{построенные супермаркеты}:
  \(1{,}6 \cdot C_{\text{spent}}\);
\item
  \(Cash\) --- деньги игрока.
\end{itemize}

\textbf{Двойной учёт.} Каждый объект входит ровно в одну составляющую,
проданное жильё --- только в \(Cash\). Если часть квартир продана в ещё
строящемся доме, её себестоимость вычитается из стоимости этого
строящегося дома.

\textbf{Выручка последнего месяца.} Выручка за месяц \(M\) обязательно
зачисляется в \(Cash\) перед подсчётом итогов. Капитал включает деньги
от всех продаж.

\section{7.
Игроки-стратегии}\label{ux438ux433ux440ux43eux43aux438-ux441ux442ux440ux430ux442ux435ux433ux438ux438}

Игроки реализуют один общий базовый класс и различаются логикой решений:

\begin{longtable}[]{@{}
  >{\raggedright\arraybackslash}p{(\columnwidth - 2\tabcolsep) * \real{0.5000}}
  >{\raggedright\arraybackslash}p{(\columnwidth - 2\tabcolsep) * \real{0.5000}}@{}}
\toprule\noalign{}
\begin{minipage}[b]{\linewidth}\raggedright
Игрок
\end{minipage} & \begin{minipage}[b]{\linewidth}\raggedright
Идея
\end{minipage} \\
\midrule\noalign{}
\endhead
\bottomrule\noalign{}
\endlastfoot
\texttt{BuilderPlayer} («застройщик») & основные деньги вкладывает в
дома, цену держит близко к справедливой, чтобы продавать первым \\
\texttt{MarketPlayer} («торговец») & строит много супермаркетов,
получает рост стоимости ×1,6, тратит на рекламу товаров \\
\texttt{BalancedPlayer} («сбалансированный») & сочетает дома и
супермаркеты, вложения делает по мере поступления выручки \\
\texttt{HumanPlayer} & решения вводит человек \\
\end{longtable}

Перед стартом для каждого места выбирается стратегия, у разных игроков
они могут быть разными (в том числе одинаковыми). Новую стратегию можно
добавить, не меняя остальные модули.

\section{8. Опциональное расширение: задержки
строительства}\label{ux43eux43fux446ux438ux43eux43dux430ux43bux44cux43dux43eux435-ux440ux430ux441ux448ux438ux440ux435ux43dux438ux435-ux437ux430ux434ux435ux440ux436ux43aux438-ux441ux442ux440ux43eux438ux442ux435ux43bux44cux441ux442ux432ux430}

Расширение модели, по умолчанию выключено; в базовой математической
модели (\texttt{math-model.md}) не используется. Если включено в
\texttt{Config}, \texttt{GameEngine} при обновлении стройки с
вероятностью \(p\) не продвигает объект на месяц (платёж за этот месяц
не списывается, срок сдвигается на месяц). Остальные модули не меняются.

\section{9. Допущения и
упрощения}\label{ux434ux43eux43fux443ux449ux435ux43dux438ux44f-ux438-ux443ux43fux440ux43eux449ux435ux43dux438ux44f}

\begin{itemize}
\tightlist
\item
  спрос и предложение измеряются в \textbf{квартирах}; деньги получаются
  умножением на площадь квартиры и цену за кв. м;
\item
  спрос считается отдельно для каждого типа дома;
\item
  реклама жилья увеличивает проданный объём игрока и не отнимает продажи
  у других;
\item
  рекламный бюджет супермаркетов --- общий на все супермаркеты игрока,
  надбавка применяется к каждому из них;
\item
  базовая прибыль супермаркета --- фиксированная сумма на один
  супермаркет;
\item
  сезонность определяется календарным месяцем, игра начинается с марта;
\item
  новая стройка начинает списывать платежи со следующего месяца после
  решения;
\item
  нехватка денег на платёж замораживает стройку на месяц;
\item
  все решения месяца принимаются одновременно;
\item
  все числовые коэффициенты --- модель по умолчанию из
  \texttt{math-model.md} (раздел 14), не заданная в ТЗ.
\end{itemize}

\section{10. Структура
проекта}\label{ux441ux442ux440ux443ux43aux442ux443ux440ux430-ux43fux440ux43eux435ux43aux442ux430}

\begin{Shaded}
\begin{Highlighting}[]
\NormalTok{project/}
\NormalTok{  config/          настройки и значения по умолчанию}
\NormalTok{  engine/          GameEngine, GameState, PublicView}
\NormalTok{  rules/           RulesValidator}
\NormalTok{  agency/          RealtorAgency}
\NormalTok{  model/           Building, House, Supermarket, Offer, Decision, District}
\NormalTok{  players/         Player, BuilderPlayer, MarketPlayer, BalancedPlayer, HumanPlayer}
\NormalTok{  ui/              интерфейс}
\NormalTok{  tests/           проверки}
\NormalTok{  docs/}
\NormalTok{    idea.md, task.md, math{-}model.md, architecture.md, images.md}
\NormalTok{    images/        готовые картинки схем}
\end{Highlighting}
\end{Shaded}

\section{11. Проверка
работы}\label{ux43fux440ux43eux432ux435ux440ux43aux430-ux440ux430ux431ux43eux442ux44b}

\begin{itemize}
\tightlist
\item
  полная партия из программных игроков для разных \(N\) и \(M\): игра
  заканчивается, победитель определён;
\item
  агентство на примере из \texttt{math-model.md} (спрос 100, предложения
  60 и 70 → продажи 60 и 40);
\item
  правила: подать неверные решения (нет денег, слишком большой объём,
  цена выше предельной, реклама без супермаркета) и убедиться, что
  \texttt{RulesValidator} их отклоняет;
\item
  начальная фаза: стартовые стройки создаются, платежи списываются в
  месяце 1;
\item
  заморозка: при нехватке денег стройка не продвигается;
\item
  капитал: проданное жильё не учитывается дважды, выручка последнего
  месяца учтена;
\item
  результаты разных стратегий (для оценки баланса можно использовать
  скрипт \texttt{balance\_sim.py}, если он включён в проект).
\end{itemize}

\end{document}
